\section*{Hoofdstuk \ref{ch:g10:storms-and-cyclones} --- Onweer en tropische cyclonen}

\begin{solution}{exo:g10:storms-and-cyclones:1}
Een \omterm{def:g10:storms-and-cyclones:thunderstorm}{onweersbui} is een bui die door een cumulonimbus wordt veroorzaakt, met
\omterm{def:g10:storms-and-cyclones:lightning}{bliksem} en donder, meestal met zware regenbuien en windstoten, soms met hagel.
\omterm{def:g10:storms-and-cyclones:lightning}{Bliksem} is een reusachtige elektrische vonk die geladen gebieden van het onweer
ontlaadt. Donder is het geluid van de lucht in het bliksemkanaal, die zo
plotseling wordt verhit dat ze met een schok uitzet.
\end{solution}

\begin{solution}{exo:g10:storms-and-cyclones:2}
Een \omterm{def:g10:storms-and-cyclones:tropical-cyclone}{tropische cycloon} is een storm van honderden kilometers breed, ontstaan
boven een warme tropische zee rond een centrum van lage druk, met winden die in
een spiraal naar binnen draaien en banden van \omterm{def:g10:storms-and-cyclones:thunderstorm}{onweersbuien}. Het \omterm{def:g10:storms-and-cyclones:tropical-cyclone}{oog} is zijn
rustige, bijna wolkeloze centrum; de \omterm{def:g10:storms-and-cyclones:tropical-cyclone}{oogwand} is de ring van torenhoge
\omterm{def:g10:storms-and-cyclones:thunderstorm}{onweersbuien} rond het \omterm{def:g10:storms-and-cyclones:tropical-cyclone}{oog}. De wind is het sterkst in de \omterm{def:g10:storms-and-cyclones:tropical-cyclone}{oogwand}.
\end{solution}

\begin{solution}{exo:g10:storms-and-cyclones:3}
\num{70}~knopen: categorie 1. \num{120}~knopen: categorie 4.
\num{140}~knopen: categorie 5. De laatste twee zijn zware orkanen (categorie 3
of hoger).
\end{solution}

\begin{solution}{exo:g10:storms-and-cyclones:4}
$343 \times 9 \approx \qty{3.1}{km}$, daarna $343 \times 15 \approx
\qty{5.1}{km}$. De tweede flits is verder weg: het onweer trekt weg.
\end{solution}

\begin{solution}{exo:g10:storms-and-cyclones:5}
\omterm{def:g10:storms-and-cyclones:lightning}{Bliksem} slaat bij voorkeur in op het hoogste punt in de omgeving: een boom die
alleen in een veld staat, is dat punt, en de stroom van een inslag loopt langs
de boom omlaag en verspreidt zich in de grond eromheen. Ga een gebouw of een
gesloten auto in en blijf daar tot ruim na de laatste donderslag.
\end{solution}

\begin{solution}{exo:g10:storms-and-cyclones:6}
Eén knoop is $1852/3600 \approx \qty{0.514}{m/s}$. \num{64}~knopen: $64 \times
0.514 \approx \qty{32.9}{m/s}$, en $32.9 \times 3.6 \approx
\qty{119}{km/h}$. \num{137}~knopen: $\qty{70.5}{m/s}$, ongeveer
\qty{254}{km/h} (de officiële tabel, via een andere eenheid afgerond, geeft
\qty{252}{km/h}).
\end{solution}

\begin{solution}{exo:g10:storms-and-cyclones:7}
$44 \times 86\,400 \approx \num{3.8e6}$ flitsen per dag, en $3.8\times10^6
\times 365.25 \approx \num{1.4e9}$, ongeveer 1.4 miljard per jaar.
\end{solution}

\begin{solution}{exo:g10:storms-and-cyclones:8}
$\Delta p = 1010 - 935 = \qty{75}{hPa}$: een bult van ongeveer \qty{75}{cm}.
\end{solution}

\begin{solution}{exo:g10:storms-and-cyclones:9}
Ze ontstaan boven de warme tropische oceanen, vooral tussen \qty{10}{\degree} en
\qty{30}{\degree} breedte op beide halfronden, met heel weinig in de zuidelijke
Atlantische Oceaan. Ze drijven naar het westen en buigen dan af naar de polen,
waar ze boven kouder water verzwakken. Op de evenaar ontbreken ze omdat de
draaiing van de Aarde de instromende lucht daar geen draai geeft: zonder draai
stroomt de lucht recht naar de lage druk en vult die op, en kan er geen
draaiende storm op gang komen.
\end{solution}

\begin{solution}{exo:g10:storms-and-cyclones:10}
$1031 - 1.56 \times 70 \approx \qty{922}{hPa}$; $1031 - 1.56 \times 40
\approx \qty{969}{hPa}$.
\end{solution}

\begin{solution}{exo:g10:storms-and-cyclones:11}
Boven land verliest de storm zijn aanvoer van verdampend water, en daarmee de
\omterm{def:g10:earth-energy-budget:latent-heat}{latente warmte} die hem aandrijft; het ruwe oppervlak remt bovendien de wind af.
Boven koud water is de verdamping veel zwakker, omdat koude lucht bij
verzadiging veel minder waterdamp bevat: de aanvoer van \omterm{def:g10:earth-energy-budget:latent-heat}{latente warmte} slinkt.
\end{solution}

\begin{solution}{exo:g10:storms-and-cyclones:12}
Bij \qty{0}{\degreeCelsius}: $\sqrt{1.4 \times 8.314 \times 273.15/0.02897}
\approx \qty{331}{m/s}$; bij \qty{30}{\degreeCelsius}, met $T =
\qty{303.15}{K}$: \qty{349}{m/s}. Voor \qty{9}{s}: $9 \times 331 \approx
\qty{2980}{m}$ en $9 \times 349 \approx \qty{3140}{m}$, een verschil van
ongeveer \qty{160}{m}, ongeveer \qty{5}{\percent}. De regel geeft \qty{3}{km}:
goed genoeg om te weten hoe dichtbij het onweer is, en meer hoeft hij niet.
\end{solution}

\begin{solution}{exo:g10:storms-and-cyclones:13}
$\Delta p = 1010 - 940 = \qty{70}{hPa}$: de bult is ongeveer \qty{0.7}{m}, het
deel van de wind $4.0 - 0.7 = \qty{3.3}{m}$, meer dan vier keer zoveel. Bij
hoogwater komt het water $4.0 + 1.5 = \qty{5.5}{m}$ boven gemiddeld zeeniveau.
\end{solution}

\begin{solution}{exo:g10:storms-and-cyclones:14}
\omterm{def:g10:storms-and-cyclones:tornado}{Tornado}: één omloop is $2\pi \times 100 \approx \qty{628}{m}$, afgelegd in
$628/90 \approx \qty{7}{s}$. Orkaan: $2\pi \times 30\,000 \approx
\qty{188}{km}$, afgelegd in $188\,000/60 \approx \qty{3140}{s}$, ongeveer
\qty{52}{min}. De \omterm{def:g10:storms-and-cyclones:tornado}{tornado} draait zo'n 450 keer sneller: hij is een kleine, heel
snelle wervel, de orkaan een enorme en tragere.
\end{solution}

\begin{solution}{exo:g10:storms-and-cyclones:15}
$24/3888 \approx \qty{0.6}{\percent}$; $414/3888 \approx
\qty{10.6}{\percent}$; $121/3888 \approx \qty{3.1}{\percent}$. Als de breedte er
niet toe deed, zou ongeveer \qty{9}{\percent} van de stormen, zo'n 350, binnen
\qty{5}{\degree} van de evenaar ontstaan; het zijn er maar 24, ongeveer vijftien
keer minder. Dicht bij de evenaar is de draaiing van de Aarde te zwak om de
storm aan het draaien te brengen.
\end{solution}

\begin{solution}{pb:g10:storms-and-cyclones:1}
\textbf{1.} $\pi \times (665\,000)^2 \approx \qty{1.39e12}{m^2}$.

\textbf{2.} $0.015 \times 1.39\times10^{12} \approx \qty{2.08e10}{m^3}$.

\textbf{3.} Eén kubieke kilometer is \qty{e9}{m^3}: ongeveer \qty{21}{km^3}.

\textbf{4.} $2.08\times10^{10} \times 1000 \approx \qty{2.08e13}{kg}$.

\textbf{5.} $2.08\times10^{13} \times 2.5\times10^6 \approx \qty{5.2e19}{J}$.

\textbf{6.} $5.2\times10^{19}/86\,400 \approx \qty{6.0e14}{W}$.

\textbf{7.} $6.0\times10^{14}/1.22\times10^{17} \approx 0.005$, ongeveer
\qty{0.5}{\percent}.

\textbf{8.} $240 \times 1.39\times10^{12} \approx \qty{3.3e14}{W}$. De storm
maakt bijna twee keer zoveel warmte vrij als het zonlicht dat zijn hele
oppervlakte absorbeert.

\textbf{9.} De storm draait niet op het zonlicht dat op hem valt: hij maakt
warmte vrij die als waterdamp is opgeslagen, verdampt uit een groter gebied van
warme zee en door de winden aangevoerd. De storm concentreert energie die over
een veel groter gebied is verzameld.

\textbf{10.} Dezelfde massa als de regen, ongeveer \qty{2.1e13}{kg} per dag (de
aangevoerde damp moet even snel worden aangevuld als hij uitregent).

\textbf{11.} $6.0\times10^{14}/82 \approx \qty{7.3e12}{m^2}$, ongeveer vijf keer
de oppervlakte van de storm: de storm moet damp uit een veel groter gebied
aanvoeren, en zijn winden laten de zee eronder veel sneller verdampen dan
gemiddeld.

\textbf{12.} \qty{24.1}{g/kg} bij \qty{28}{\degreeCelsius} tegen
\qty{14.7}{g/kg} bij \qty{20}{\degreeCelsius}: $24.1/14.7 \approx 1.64$,
ongeveer \qty{64}{\percent} meer.

\textbf{13.} Warm water verdampt snel en warme lucht bevat veel meer damp: de
instromende lucht voert veel meer \omterm{def:g10:earth-energy-budget:latent-heat}{latente warmte} aan om vrij te maken, en
verzadigd als ze is, stijgt ze vrij op. Boven koelere zeeën is de aanvoer te
klein om de motor aan te drijven.

\textbf{14.} Boven land is er geen zee om te verdampen: de aanvoer van \omterm{def:g10:earth-energy-budget:latent-heat}{latente
warmte} stopt, en de ruwe grond remt de wind af.

\textbf{15.} $1.5\times10^{12} \times 86\,400 \approx \qty{1.3e17}{J}$.

\textbf{16.} $96 \times 1852/3600 \approx \qty{49.4}{m/s}$.

\textbf{17.} $\frac12 \times 1.16 \times 49.4^2 \approx \qty{1.4e3}{J}$.

\textbf{18.} $5.2\times10^{19}/1.3\times10^{17} = 400$: de \omterm{def:g10:earth-energy-budget:latent-heat}{latente warmte} is
zo'n 400 keer de energie die aan de winden wordt gegeven.

\textbf{19.} Het grootste deel verwarmt de bovenlucht van de storm, die zich
bovenaan uitspreidt en de warmte als infraroodstraling aan de ruimte afgeeft;
de winden zelf verliezen hun energie door wrijving met de zee en de grond.

\textbf{20.} $1/400 = \qty{0.25}{\percent}$ (het script geeft een verhouding van
402): \textbf{ongeveer één deel op 400 van de warmte die bij condensatie
vrijkomt, wordt wind}.
\end{solution}
